\chapter{Závěr}
Cílem této práce bylo navrhnout a~implementovat informační systém pro automatizaci provozu turistického kempu Olšovec.
V~první kapitole jsem prostudoval legislativní rámec pro provozovatele krátkodobých ubytovacích služeb, který provozovatel kempu ocenil jako ucelený přehled aktuálního právního stavu.
Následně jsem se zabýval problematikou autentizace a~autorizace a~bezpečností webových aplikací; získaná zjištění byla zohledněna v~průběhu implementace.

Před samotným návrhem byla provedena rešerše existujících řešení, aby nedocházelo k~vývoji již dostupné funkcionality.
Procesy rezervací, evidence a~odbavení hostů byly analyzovány ve~spolupráci s~Kempem Olšovec a~na~jejich základě byly definovány požadavky na~automatizaci.
Návrh systému kladl důraz na~bezpečnost, použitelnost, oddělení jednotlivých vrstev, internacionalizaci; řešeno bylo rovněž napojení na~automatické otevírání vjezdových závor pro průjezd vozidel.
Navržený systém byl následně implementován a~ve~spolupráci s~kempem otestován; součástí výstupu je rovněž dokumentace pro jednoduché zprovoznění.
Největší komplikací byla nedostatečná dokumentace systému Ubyport.

Výsledný systém je realizován jako progresivní webová aplikace pomocí 4 klientských variant určených pro recepční, hosty u~recepčního pultu, provozní personál v~areálu a~veřejnost rezervující si pobyt.
Veřejný rezervační portál nabízí vícekrokového průvodce rezervací a~online check-in.
Tablet umístěný na~recepčním pultu slouží k~pořízení digitálního podpisu hosta pro kompletní omezení užití papíru.
Implementovaný systém umožňuje napojení na~vjezdový systém závor, pro časovou úsporu při odbavování hostů.
Systém je plně funkční, otestovaný a~připravený k~nasazení do~produkce; v~současné době je využívána jeho neveřejná část.
